\let\negmedspace\undefined
\let\negthickspace\undefined
\documentclass[journal]{article}
\usepackage[a5paper, margin=10mm, onecolumn]{geometry}
\usepackage{lmodern} % Ensure lmodern is loaded for pdflatex

\setlength{\headheight}{1cm} % Set the height of the header box
\setlength{\headsep}{0mm}     % Set the distance between the header box and the top of the text

\usepackage{gvv-book}
\usepackage{gvv}
\usepackage{cite}
\usepackage{textcomp}
\usepackage{amsmath,amssymb,amsfonts,amsthm}
\usepackage{algorithmic}
\usepackage{graphicx}
\graphicspath{{./figs/}}
\usepackage{textcomp}
\usepackage{xcolor}
\usepackage{txfonts}
\usepackage{listings}
\usepackage{enumitem}
\usepackage{mathtools}
\usepackage{gensymb}
\usepackage{comment}
\usepackage[breaklinks=true]{hyperref}
\usepackage{tkz-euclide} 
\usepackage{listings}
\usepackage{gvv}                                        
\def\inputGnumericTable{}                                 
\usepackage[latin1]{inputenc}                                
\usepackage{color}                                            
\usepackage{array}                                            
\usepackage{longtable}                                       
\usepackage{calc}                                             
\usepackage{multirow}                                         
\usepackage{hhline}                                           
\usepackage{ifthen}                                           
\usepackage{lscape}
\usepackage{circuitikz}
\tikzstyle{block} = [rectangle, draw, fill=blue!20, 
text width=4em, text centered, rounded corners, minimum height=3em]
\tikzstyle{sum} = [draw, fill=blue!10, circle, minimum size=1cm, node distance=1.5cm]
\tikzstyle{input} = [coordinate]
\tikzstyle{output} = [coordinate]


\begin{document}
	
	\bibliographystyle{IEEEtran}
	\vspace{3cm}
	
\title{12.164}
\author{EE25BTECH11047 - RAVULA SHASHANK REDDY}
\maketitle
\hrulefill
\bigskip 

\renewcommand{\thetable}{\theenumi}
\setlength{\intextsep}{10pt}

\textbf{Question:} \\

Let 
\begin{align*}
\vec{A} = \myvec{1 & 1 & 1 \\ 0 & 1 & 1 \\ 0 & 0 & 1}.
\end{align*}

Find the maximum number of linearly independent eigenvectors of \(\vec{A}\).

\textbf{Solution:}\\

Since $\vec{A}$ is a triangular matrix
\begin{align}
 & \lambda = 1,1,1 \\
(\vec{A} - \vec{I})\vec{x} &= \myvec{0 & 1 & 1 \\ 0 & 0 & 1 \\ 0 & 0 & 0} \vec{x} = 0, \quad \vec{x} = \myvec{x \\ y \\ z} \\
& y + z = 0 \implies y = -z \\
& z = 0 \implies y = 0 \\
\vec{x} &= \myvec{x \\ 0 \\ 0} = x \myvec{1 \\ 0 \\ 0}
\end{align}


Maximum number of linearly independent eigenvectors = 1


\end{document}